% ============================================================
% Question Bank: ch08-03 Comparing Two Populations Visually
% Topic Title: Comparing Two Populations Visually
% CCSS: 7.SP.B.3
% Questions: 27 (15 MC + 6 SA + 6 Graphical)
% ============================================================

% --- Multiple Choice Questions (15) ---

\begin{practiceQuestion}{8-3-q01}{mc}
\begin{questionText}
Two soccer teams recorded their goals per game. Team A has a mean of $3.2$ goals with a MAD of $0.8$. Team B has a mean of $2.4$ goals with a MAD of $0.8$. The difference in their means expressed as a multiple of the MAD is:
\end{questionText}
\choiceA{$0.8$}
\choiceB{$1.0$}
\choiceC{$1.6$}
\choiceD{$4.0$}
\correctAnswer{B}
\explanation{Difference in means: $3.2 - 2.4 = 0.8$. Divide by the MAD: $\frac{0.8}{0.8} = 1.0$ MAD. A difference of about $1$ MAD suggests some separation but substantial overlap.}
\end{practiceQuestion}

\begin{practiceQuestion}{8-3-q02}{mc}
\begin{questionText}
If the difference between two data set means is less than $1$ MAD, you can conclude that:
\end{questionText}
\choiceA{The two populations are definitely the same}
\choiceB{The distributions overlap significantly and the difference is not meaningful}
\choiceC{One group is clearly better than the other}
\choiceD{The MAD was calculated incorrectly}
\correctAnswer{B}
\explanation{When the difference in means is less than $1$ MAD, the data sets have significant overlap, suggesting the difference between groups is not meaningful.}
\end{practiceQuestion}

\begin{practiceQuestion}{8-3-q03}{mc}
\begin{questionText}
Which display is most useful for comparing the shapes of two small data sets side by side?
\end{questionText}
\choiceA{A single circle graph}
\choiceB{Side-by-side dot plots on the same number line}
\choiceC{A frequency table with one column}
\choiceD{A single bar graph}
\correctAnswer{B}
\explanation{Side-by-side dot plots on the same number line let you see the center, spread, and overlap of two distributions at a glance.}
\end{practiceQuestion}

\begin{practiceQuestion}{8-3-q04}{mc}
\begin{questionText}
Group P: median $45$, IQR $6$. Group Q: median $52$, IQR $6$. How many IQRs apart are the medians?
\end{questionText}
\choiceA{$\frac{7}{6} \approx 1.2$ IQRs}
\choiceB{$6$ IQRs}
\choiceC{$7$ IQRs}
\choiceD{$0.5$ IQRs}
\correctAnswer{A}
\explanation{Difference in medians: $52 - 45 = 7$. Divide by the IQR: $\frac{7}{6} \approx 1.17$ IQRs. This suggests some separation between the groups but still overlap.}
\end{practiceQuestion}

\begin{practiceQuestion}{8-3-q05}{mc}
\begin{questionText}
Two box plots are placed on the same number line. The boxes do not overlap at all. What can you conclude?
\end{questionText}
\choiceA{The two groups performed about the same}
\choiceB{The middle $50\%$ of each group's data does not share any values}
\choiceC{Both groups have the same median}
\choiceD{Both groups have the same range}
\correctAnswer{B}
\explanation{The box in a box plot represents the middle $50\%$ of the data (from $Q_1$ to $Q_3$). If the boxes do not overlap, the central halves of the two distributions have no values in common.}
\end{practiceQuestion}

\begin{practiceQuestion}{8-3-q06}{mc}
\begin{questionText}
Class X has a mean score of $88$ with a MAD of $4$. Class Y has a mean score of $80$ with a MAD of $4$. What is the difference in means as a multiple of the MAD?
\end{questionText}
\choiceA{$0.5$ MADs}
\choiceB{$1$ MAD}
\choiceC{$2$ MADs}
\choiceD{$4$ MADs}
\correctAnswer{C}
\explanation{Difference: $88 - 80 = 8$. Divide by the MAD: $\frac{8}{4} = 2$ MADs. A difference of $2$ or more MADs typically indicates a meaningful difference between the groups.}
\end{practiceQuestion}

\begin{practiceQuestion}{8-3-q07}{mc}
\begin{questionText}
Which of the following is true about two histograms that have a large amount of overlap?
\end{questionText}
\choiceA{The populations are very different}
\choiceB{The populations are similar in terms of center and spread}
\choiceC{One histogram must be shifted far to the right}
\choiceD{The standard deviations must be zero}
\correctAnswer{B}
\explanation{Large overlap in histograms means the data values in both groups are distributed similarly. This indicates the populations have similar centers and spreads.}
\end{practiceQuestion}

\begin{practiceQuestion}{8-3-q08}{mc}
\begin{questionText}
A researcher compares the running speeds (in mph) of two dog breeds. Breed A: mean $18$, MAD $2$. Breed B: mean $15$, MAD $2$. A student says the difference is not meaningful. Is the student correct?
\end{questionText}
\choiceA{Yes, because both MADs are equal}
\choiceB{No, because $\frac{3}{2} = 1.5$ MADs is a meaningful difference}
\choiceC{Yes, because the means are close to each other}
\choiceD{No, because the difference is exactly $3$ mph}
\correctAnswer{B}
\explanation{Difference in means: $18 - 15 = 3$. In terms of MAD: $\frac{3}{2} = 1.5$ MADs. A difference of $1.5$ MADs or more generally suggests a meaningful difference between the two groups.}
\end{practiceQuestion}

\begin{practiceQuestion}{8-3-q09}{mc}
\begin{questionText}
Two dot plots show quiz scores for two classes. Class M's dots are clustered tightly around $85$, while Class N's dots are spread from $60$ to $100$. Which statement best describes the comparison?
\end{questionText}
\choiceA{Class M and Class N have the same center}
\choiceB{Class N is more consistent than Class M}
\choiceC{Class M is more consistent than Class N}
\choiceD{Class N always outperforms Class M}
\correctAnswer{C}
\explanation{Tightly clustered data means less variability. Class M's scores are concentrated near $85$, showing more consistency. Class N's wide spread indicates more variability in performance.}
\end{practiceQuestion}

\begin{practiceQuestion}{8-3-q10}{mc}
\begin{questionText}
To say the difference between two groups is ``meaningful,'' the difference in centers should typically be at least:
\end{questionText}
\choiceA{$0.5$ MADs}
\choiceB{$1$ MAD}
\choiceC{$2$ MADs or more}
\choiceD{$10$ MADs}
\correctAnswer{C}
\explanation{A common guideline is that a difference of $2$ or more MADs between the centers suggests a meaningful difference, indicating that the populations are likely different.}
\end{practiceQuestion}

\begin{practiceQuestion}{8-3-q11}{mc}
\begin{questionText}
The box plots for two basketball teams show that Team A's median is $22$ points and Team B's median is $18$ points. Both have a MAD of $3$. Which conclusion is best supported?
\end{questionText}
\choiceA{The teams are equally skilled}
\choiceB{Team A tends to score about $\frac{4}{3} \approx 1.3$ MADs more than Team B}
\choiceC{Team B always loses}
\choiceD{The data sets do not overlap}
\correctAnswer{B}
\explanation{Difference in medians: $22 - 18 = 4$. In MADs: $\frac{4}{3} \approx 1.3$. This suggests Team A tends to score higher, with some overlap between the teams.}
\end{practiceQuestion}

\begin{practiceQuestion}{8-3-q12}{mc}
\begin{questionText}
Two dot plots show the number of books students read in Class G and Class H. Class G's distribution is centered at $5$ books and Class H's at $5$ books, but Class H's distribution is much wider. What does this tell you?
\end{questionText}
\choiceA{Both classes read the same number of books per student}
\choiceB{The classes have similar centers but different variability}
\choiceC{Class H reads more books than Class G}
\choiceD{The data cannot be compared}
\correctAnswer{B}
\explanation{Same center (median or mean around $5$) means both classes read about the same number of books on average. A wider distribution for Class H means more variability — some students read very few and others read many.}
\end{practiceQuestion}

\begin{practiceQuestion}{8-3-q13}{mc}
\begin{questionText}
Store A: mean daily revenue $\$420$, MAD $\$30$. Store B: mean daily revenue $\$480$, MAD $\$30$. The difference in means expressed in MADs is:
\end{questionText}
\choiceA{$0.5$ MADs}
\choiceB{$1$ MAD}
\choiceC{$2$ MADs}
\choiceD{$3$ MADs}
\correctAnswer{C}
\explanation{Difference: $\$480 - \$420 = \$60$. Divide by the MAD: $\frac{60}{30} = 2$ MADs. This is a meaningful difference, indicating Store B consistently earns more.}
\end{practiceQuestion}

\begin{practiceQuestion}{8-3-q14}{mc}
\begin{questionText}
A student compares two data sets using dot plots. Set 1 has $10$ values all between $20$ and $25$. Set 2 has $10$ values between $15$ and $35$. Which set has a larger MAD?
\end{questionText}
\choiceA{Set 1}
\choiceB{Set 2}
\choiceC{They have the same MAD}
\choiceD{MAD cannot be determined from this information}
\correctAnswer{B}
\explanation{The MAD measures the average distance from the mean. Set 2 has data spread over a wider range ($15$ to $35$), so its values are farther from the mean on average, giving a larger MAD.}
\end{practiceQuestion}

\begin{practiceQuestion}{8-3-q15}{mc}
\begin{questionText}
Two histograms show the ages of participants in two programs. Program A's histogram is shifted to the right compared to Program B. What does this indicate?
\end{questionText}
\choiceA{Program A has younger participants}
\choiceB{Program B has older participants}
\choiceC{Program A has older participants on average}
\choiceD{Both programs have the same age distribution}
\correctAnswer{C}
\explanation{A histogram shifted to the right indicates higher values. Program A being shifted right means its participants tend to be older on average than Program B's participants.}
\end{practiceQuestion}

% --- Short Answer Questions (6) ---

\begin{practiceQuestion}{8-3-q16}{sa}
\begin{questionText}
Group R has a mean of $36$ and a MAD of $5$. Group S has a mean of $46$. Express the difference in means as a multiple of the MAD.
\end{questionText}
\correctAnswer{$2$}
\explanation{Difference in means: $46 - 36 = 10$. Divide by the MAD: $\frac{10}{5} = 2$ MADs. This is a meaningful difference between the two groups.}
\end{practiceQuestion}

\begin{practiceQuestion}{8-3-q17}{sa}
\begin{questionText}
Team C: mean $65$, MAD $7$. Team D: mean $72$, MAD $7$. Is the difference in means meaningful? Justify using MADs.
\end{questionText}
\correctAnswer{No; the difference is $1$ MAD, so it is not clearly meaningful.}
\explanation{Difference: $72 - 65 = 7$. In MADs: $\frac{7}{7} = 1$. A difference of $1$ MAD shows only modest separation with considerable overlap, so it is not usually considered a meaningful difference.}
\end{practiceQuestion}

\begin{practiceQuestion}{8-3-q18}{sa}
\begin{questionText}
Data set A has a median of $40$ and a range of $12$. Data set B has a median of $40$ and a range of $28$. Which data set has more variability?
\end{questionText}
\correctAnswer{Data set B}
\explanation{Both sets have the same median ($40$), but Data set B has a range of $28$ compared to $12$. A larger range means the data is more spread out, so B has more variability.}
\end{practiceQuestion}

\begin{practiceQuestion}{8-3-q19}{sa}
\begin{questionText}
The test scores for School A have a MAD of $3$ and for School B a MAD of $9$. Their means are both $75$. Which school's scores are more consistent?
\end{questionText}
\correctAnswer{School A}
\explanation{MAD measures how far scores typically are from the mean. A smaller MAD ($3$ vs $9$) means scores are closer to the mean on average, so School A's scores are more consistent.}
\end{practiceQuestion}

\begin{practiceQuestion}{8-3-q20}{sa}
\begin{questionText}
Plant heights in Garden A: $12, 14, 15, 16, 18$ cm. Find the MAD of Garden A's plant heights.
\end{questionText}
\correctAnswer{$1.6$}
\explanation{Mean: $\frac{12+14+15+16+18}{5} = \frac{75}{5} = 15$. Absolute deviations: $|12-15|=3$, $|14-15|=1$, $|15-15|=0$, $|16-15|=1$, $|18-15|=3$. MAD: $\frac{3+1+0+1+3}{5} = \frac{8}{5} = 1.6$ cm.}
\end{practiceQuestion}

\begin{practiceQuestion}{8-3-q21}{sa}
\begin{questionText}
Two data sets have means that differ by $6$ units. The MAD for both sets is $4$. Express the difference in means as a multiple of the MAD and state whether the difference is likely meaningful.
\end{questionText}
\correctAnswer{$1.5$ MADs; the difference is likely meaningful.}
\explanation{Divide the difference by the MAD: $\frac{6}{4} = 1.5$ MADs. A difference of $1.5$ MADs indicates noticeable separation between the two populations, suggesting the difference is likely meaningful.}
\end{practiceQuestion}

% --- Graphical Questions (3 GMC + 3 GSA) ---

\begin{practiceQuestion}{8-3-q22}{gmc}
\begin{questionText}
The dot plots below show the number of push-ups completed by students in two PE classes.

\begin{center}
\begin{tikzpicture}[scale=0.45]
  \draw[thick,->] (7.5,0) -- (22,0);
  \foreach \x in {8,...,21} \draw (\x,0.15) -- (\x,-0.15) node[below,font=\tiny] {$\x$};
  \node[below] at (14.5,-1) {Push-ups};
  % Class 1 dots (center ~14)
  \foreach \y in {1} \fill[funBlue] (10,0.1+\y*0.4) circle (4pt);
  \foreach \y in {1,2} \fill[funBlue] (12,0.1+\y*0.4) circle (4pt);
  \foreach \y in {1,2,3} \fill[funBlue] (13,0.1+\y*0.4) circle (4pt);
  \foreach \y in {1,2,3,4} \fill[funBlue] (14,0.1+\y*0.4) circle (4pt);
  \foreach \y in {1,2,3} \fill[funBlue] (15,0.1+\y*0.4) circle (4pt);
  \foreach \y in {1,2} \fill[funBlue] (16,0.1+\y*0.4) circle (4pt);
  \foreach \y in {1} \fill[funBlue] (18,0.1+\y*0.4) circle (4pt);
  \node[funBlue,above] at (14,2.5) {\small Class 1};
  % Class 2 dots (center ~17)
  \foreach \y in {1} \fill[funRed] (13,0.1+\y*0.4-3.5) circle (4pt);
  \foreach \y in {1,2} \fill[funRed] (15,0.1+\y*0.4-3.5) circle (4pt);
  \foreach \y in {1,2,3} \fill[funRed] (16,0.1+\y*0.4-3.5) circle (4pt);
  \foreach \y in {1,2,3,4} \fill[funRed] (17,0.1+\y*0.4-3.5) circle (4pt);
  \foreach \y in {1,2,3} \fill[funRed] (18,0.1+\y*0.4-3.5) circle (4pt);
  \foreach \y in {1,2} \fill[funRed] (19,0.1+\y*0.4-3.5) circle (4pt);
  \foreach \y in {1} \fill[funRed] (21,0.1+\y*0.4-3.5) circle (4pt);
  \node[funRed,below] at (17,-3.5) {\small Class 2};
  \draw[thick,->] (7.5,-3.3) -- (22,-3.3);
  \foreach \x in {8,...,21} \draw (\x,-3.15) -- (\x,-3.45) node[below,font=\tiny] {$\x$};
\end{tikzpicture}
\end{center}

Which statement best compares the two classes?
\end{questionText}
\choiceA{Class 2's center is higher than Class 1's with partial overlap}
\choiceB{Class 1 and Class 2 have the same center}
\choiceC{Class 1's center is higher than Class 2's}
\choiceD{The two classes have no overlap at all}
\correctAnswer{A}
\explanation{Class 1 is centered around $14$ push-ups and Class 2 around $17$. The distributions overlap in the range of about $13$--$18$, but Class 2 is shifted higher on average.}
\end{practiceQuestion}

\begin{practiceQuestion}{8-3-q23}{gmc}
\begin{questionText}
The side-by-side box plots below show the daily high temperatures (°F) in two cities during June.

\begin{center}
\begin{tikzpicture}[scale=0.35]
  \draw[thick,->] (59,0) -- (101,0);
  \foreach \x in {60,65,...,100} \draw (\x,0.2) -- (\x,-0.2) node[below] {\tiny $\x$};
  \node[below] at (80,-1.2) {Temperature (°F)};
  % City A: min=65, Q1=70, med=75, Q3=82, max=90
  \draw[thick] (65,2) -- (70,2); % whisker
  \draw[thick,fill=funBlue!30] (70,1) rectangle (82,3);
  \draw[very thick,funBlue] (75,1) -- (75,3);
  \draw[thick] (82,2) -- (90,2);
  \draw[thick] (65,1.5) -- (65,2.5);
  \draw[thick] (90,1.5) -- (90,2.5);
  \node[right] at (91,2) {\small City A};
  % City B: min=72, Q1=78, med=84, Q3=90, max=98
  \draw[thick] (72,5) -- (78,5);
  \draw[thick,fill=funRed!30] (78,4) rectangle (90,6);
  \draw[very thick,funRed] (84,4) -- (84,6);
  \draw[thick] (90,5) -- (98,5);
  \draw[thick] (72,4.5) -- (72,5.5);
  \draw[thick] (98,4.5) -- (98,5.5);
  \node[right] at (99,5) {\small City B};
\end{tikzpicture}
\end{center}

What is the difference in the medians of the two cities?
\end{questionText}
\choiceA{$5°$F}
\choiceB{$7°$F}
\choiceC{$9°$F}
\choiceD{$12°$F}
\correctAnswer{C}
\explanation{City A's median is $75°$F and City B's median is $84°$F. The difference is $84 - 75 = 9°$F. City B was typically warmer.}
\end{practiceQuestion}

\begin{practiceQuestion}{8-3-q24}{gmc}
\begin{questionText}
Two histograms show weekly study hours for students in two programs.

\begin{center}
\begin{tikzpicture}[scale=0.5]
  % Program A
  \node[above] at (3,6) {\textbf{Program A}};
  \draw[thick,->] (0,0) -- (0,5.5);
  \draw[thick,->] (0,0) -- (6.5,0);
  \foreach \y in {1,...,5} {\draw (-0.15,\y) -- (0,\y); \node[left,font=\tiny] at (0,\y) {$\y$};}
  \fill[funBlue!60] (0,0) rectangle (1,1);
  \fill[funBlue!60] (1,0) rectangle (2,3);
  \fill[funBlue!60] (2,0) rectangle (3,5);
  \fill[funBlue!60] (3,0) rectangle (4,4);
  \fill[funBlue!60] (4,0) rectangle (5,2);
  \fill[funBlue!60] (5,0) rectangle (6,1);
  \node[below,font=\tiny] at (0.5,0) {$2$};
  \node[below,font=\tiny] at (1.5,0) {$4$};
  \node[below,font=\tiny] at (2.5,0) {$6$};
  \node[below,font=\tiny] at (3.5,0) {$8$};
  \node[below,font=\tiny] at (4.5,0) {$10$};
  \node[below,font=\tiny] at (5.5,0) {$12$};
  % Program B
  \begin{scope}[xshift=9cm]
  \node[above] at (3,6) {\textbf{Program B}};
  \draw[thick,->] (0,0) -- (0,5.5);
  \draw[thick,->] (0,0) -- (6.5,0);
  \foreach \y in {1,...,5} {\draw (-0.15,\y) -- (0,\y); \node[left,font=\tiny] at (0,\y) {$\y$};}
  \fill[funRed!60] (0,0) rectangle (1,2);
  \fill[funRed!60] (1,0) rectangle (2,4);
  \fill[funRed!60] (2,0) rectangle (3,5);
  \fill[funRed!60] (3,0) rectangle (4,3);
  \fill[funRed!60] (4,0) rectangle (5,1);
  \fill[funRed!60] (5,0) rectangle (6,1);
  \node[below,font=\tiny] at (0.5,0) {$4$};
  \node[below,font=\tiny] at (1.5,0) {$6$};
  \node[below,font=\tiny] at (2.5,0) {$8$};
  \node[below,font=\tiny] at (3.5,0) {$10$};
  \node[below,font=\tiny] at (4.5,0) {$12$};
  \node[below,font=\tiny] at (5.5,0) {$14$};
  \end{scope}
\end{tikzpicture}
\end{center}

Which statement best compares the two programs?
\end{questionText}
\choiceA{Program B's distribution is shifted to the right compared to Program A}
\choiceB{Both programs have the same distribution}
\choiceC{Program A's distribution is shifted to the right}
\choiceD{Program A has more variability}
\correctAnswer{A}
\explanation{Program A's bars span $2$--$12$ hours while Program B's span $4$--$14$ hours. Program B's histogram is shifted to the right, meaning Program B students tend to study more hours per week.}
\end{practiceQuestion}

\begin{practiceQuestion}{8-3-q25}{gsa}
\begin{questionText}
The dot plots below show daily text messages sent by students in two friend groups.

\begin{center}
\begin{tikzpicture}[scale=0.5]
  \draw[thick,->] (-0.5,0) -- (12,0);
  \foreach \x in {0,...,11} \draw (\x,0.15) -- (\x,-0.15) node[below,font=\tiny] {$\x$};
  \node[below] at (5.5,-0.9) {Text Messages};
  % Group A: 2,3,3,4,4,4,5,5,6 → mean = 36/9 = 4
  \fill[funBlue] (2,0.5) circle (4pt);
  \foreach \y in {1,2} \fill[funBlue] (3,0.1+\y*0.4) circle (4pt);
  \foreach \y in {1,2,3} \fill[funBlue] (4,0.1+\y*0.4) circle (4pt);
  \foreach \y in {1,2} \fill[funBlue] (5,0.1+\y*0.4) circle (4pt);
  \fill[funBlue] (6,0.5) circle (4pt);
  \node[funBlue,above] at (4,2) {\small Group A};
  % Group B: 5,6,6,7,7,7,8,8,9 → mean = 63/9 = 7
  \fill[funRed] (5,0.5+3) circle (4pt);
  \foreach \y in {1,2} \fill[funRed] (6,0.1+\y*0.4+3) circle (4pt);
  \foreach \y in {1,2,3} \fill[funRed] (7,0.1+\y*0.4+3) circle (4pt);
  \foreach \y in {1,2} \fill[funRed] (8,0.1+\y*0.4+3) circle (4pt);
  \fill[funRed] (9,0.5+3) circle (4pt);
  \node[funRed,above] at (7,5) {\small Group B};
\end{tikzpicture}
\end{center}

Find the mean for each group and the difference in means.
\end{questionText}
\correctAnswer{Group A mean $= 4$; Group B mean $= 7$; difference $= 3$}
\explanation{Group A: $\frac{2+3+3+4+4+4+5+5+6}{9} = \frac{36}{9} = 4$. Group B: $\frac{5+6+6+7+7+7+8+8+9}{9} = \frac{63}{9} = 7$. Difference: $7 - 4 = 3$ messages.}
\end{practiceQuestion}

\begin{practiceQuestion}{8-3-q26}{gsa}
\begin{questionText}
The box plots below compare the weights (in ounces) of apples from two orchards.

\begin{center}
\begin{tikzpicture}[scale=0.5]
  \draw[thick,->] (3,0) -- (13,0);
  \foreach \x in {4,...,12} \draw (\x,0.2) -- (\x,-0.2) node[below,font=\small] {$\x$};
  \node[below] at (8,-1) {Weight (oz)};
  % Orchard A: min=5, Q1=6, med=7, Q3=8, max=10
  \draw[thick] (5,2) -- (6,2);
  \draw[thick,fill=funBlue!30] (6,1.3) rectangle (8,2.7);
  \draw[very thick,funBlue] (7,1.3) -- (7,2.7);
  \draw[thick] (8,2) -- (10,2);
  \draw[thick] (5,1.6) -- (5,2.4);
  \draw[thick] (10,1.6) -- (10,2.4);
  \node[left] at (4.5,2) {\small Orchard A};
  % Orchard B: min=6, Q1=7, med=9, Q3=10, max=12
  \draw[thick] (6,4.5) -- (7,4.5);
  \draw[thick,fill=funRed!30] (7,3.8) rectangle (10,5.2);
  \draw[very thick,funRed] (9,3.8) -- (9,5.2);
  \draw[thick] (10,4.5) -- (12,4.5);
  \draw[thick] (6,4.1) -- (6,4.9);
  \draw[thick] (12,4.1) -- (12,4.9);
  \node[left] at (5.5,4.5) {\small Orchard B};
\end{tikzpicture}
\end{center}

Find the IQR for each orchard and state which orchard produces heavier apples on average.
\end{questionText}
\correctAnswer{Orchard A IQR $= 2$ oz; Orchard B IQR $= 3$ oz; Orchard B is heavier.}
\explanation{Orchard A: IQR $= Q_3 - Q_1 = 8 - 6 = 2$ oz, median $= 7$ oz. Orchard B: IQR $= 10 - 7 = 3$ oz, median $= 9$ oz. Orchard B has a higher median, so it produces heavier apples on average.}
\end{practiceQuestion}

\begin{practiceQuestion}{8-3-q27}{gsa}
\begin{questionText}
The dot plots show hours of community service completed by members of two clubs.

\begin{center}
\begin{tikzpicture}[scale=0.5]
  % Club A number line
  \draw[thick,->] (3.5,0) -- (14,0);
  \foreach \x in {4,...,13} \draw (\x,0.15) -- (\x,-0.15) node[below,font=\tiny] {$\x$};
  % Club A: 6,7,7,8,8,8,9,9,10 sum=72 mean=8
  \fill[funBlue] (6,0.5) circle (4pt);
  \foreach \y in {1,2} \fill[funBlue] (7,0.1+\y*0.4) circle (4pt);
  \foreach \y in {1,2,3} \fill[funBlue] (8,0.1+\y*0.4) circle (4pt);
  \foreach \y in {1,2} \fill[funBlue] (9,0.1+\y*0.4) circle (4pt);
  \fill[funBlue] (10,0.5) circle (4pt);
  \node[funBlue,above] at (8,1.6) {\small Club A};

  \begin{scope}[yshift=-4cm]
    % Club B number line
    \draw[thick,->] (3.5,0) -- (14,0);
    \foreach \x in {4,...,13} \draw (\x,0.15) -- (\x,-0.15) node[below,font=\tiny] {$\x$};
    % Club B: 9,10,10,11,11,11,12,12,13 sum=99 mean=11
    \fill[funRed] (9,0.5) circle (4pt);
    \foreach \y in {1,2} \fill[funRed] (10,0.1+\y*0.4) circle (4pt);
    \foreach \y in {1,2,3} \fill[funRed] (11,0.1+\y*0.4) circle (4pt);
    \foreach \y in {1,2} \fill[funRed] (12,0.1+\y*0.4) circle (4pt);
    \fill[funRed] (13,0.5) circle (4pt);
    \node[funRed,above] at (11,1.6) {\small Club B};
  \end{scope}

  \node[below] at (8.5,-5.2) {Hours};
\end{tikzpicture}
\end{center}

Find the mean for each club. Then find the MAD for Club A and express the difference in means as a multiple of the MAD.
\end{questionText}
\correctAnswer{Club A mean $= 8$, Club B mean $= 11$; MAD of Club A $\approx 0.9$; difference $\approx 3.4$ MADs}
\explanation{Club A: $\frac{6+7+7+8+8+8+9+9+10}{9} = \frac{72}{9} = 8$. Deviations from $8$: $2,1,1,0,0,0,1,1,2$; MAD $= \frac{8}{9} \approx 0.89$. Club B mean: $\frac{9+10+10+11+11+11+12+12+13}{9} = \frac{99}{9} = 11$. Difference: $\frac{11-8}{0.89} \approx 3.4$ MADs, indicating a meaningful difference.}
\end{practiceQuestion}
